\scanpage{56}
\mspage{55}

c) $\varnothing\in\unsure{\Phi|Q|\varphi}$ \deleted{\unsure{$\ne\varnothing$} [en vertu de b), c'est tjrs vrai si $\unsure{\mathcal M|\mathscr M}\ne\varnothing$], et \unsure{si} $\mathcal M=\varnothing$, \unsure{cela signifie que} \unsure{$\varnothing$}}

\deleted{\uline{\unsure{Rép.|Rq.}} (équivaut : $\varnothing\in\Phi$)}

\uline{Remarque} Inversement, soient $\mathcal M$ un ens, $\deleted{\Phi}\unsure{\mathcal F|\mathscr F}\subset\unsure{\mathcal P|\mathscr P}(\mathcal M)$ satisfaisant les conditions a) b) c). NB La condition c résulte de \inserted{a),} b si $\mathcal M\ne\varnothing$. Dans tous les cas, elle équivaut à : \deleted{\unsure{$\in$|à}} c') $\deleted{\Phi}\mathcal F\ne\varnothing$ \unsure{i.e} existe $\deleted{X}A\subset\mathcal M$ telle que $A\in\deleted{\Phi}\mathcal F$

% slanted
\margin{\unsure{Ces|Les} conditions impliquent que $\mathcal F$ est stable par $A\cap B$ (mais en général on n'a \unsure{$\not\Rightarrow$|pas} $\mathcal M\in\mathcal F$ !)}

Alors

1°) $\Phi$ ordonné par inclusion \deleted{\unsure{ad|sa}}satisfait les conditions C1, C2. \deleted{Les objets irréductibles} \deleted{\unsure{sont}}

2°) Les objets irréductibles sont les $A_X$, et $X\mapsto\deleted{\unsure{X}}A_X$ est une bijection de $\mathcal M$ \deleted{avec} l'ens des objets \unsure{irréd.} de $\deleted{\Phi}\mathcal F$.

3°) Si $X\in\mathcal M$, $A\in\deleted{\Phi}\mathcal F$, alors $X\in A\Longleftrightarrow A_X\subset A$ \deleted{, donc} (tautologique) donc \deleted{\unsure{$A\in\mathcal F$}} si on identifie $\mathcal M$ \inserted{(par $X\unsure{\vdash|\mapsto} A_X$)} à l'ens des objets irréductibles de $\deleted{\Phi}\mathcal F$, alors pour $A\in\mathcal F$, $\overline A=\{A_X\mid A_X\subset A\inserted{\text{\unsure{ie.}}\,X\in A}\}$ s'identifie à $A$ lui-$=$.

\centerline{\rule{1.8cm}{0.35pt}}

On \unsure{considérera|considère} $\mathcal M$ comme ordonné par $\leq$ (jamais par $\unsure{\prec|\preceq}$ !) : multistrates et \uline{sous-multistrates}. Les figures \inserted{$F$} peuvent s'identifier aux \deleted{\unsure{des parties}} ens. de multistrates $\overline F$ correspondants (chaque figure à l'ens de ses strates).
